\documentclass{article}
\usepackage{amsmath,amssymb}
\usepackage[vocab={curves}]{leanify}
\newtheorem{lemma}{Lemma}
\begin{document}

\section{Notation}

We write $v(x,y,t)$ for the speed of the curve with coordinates $x$ and $y$ at $t$.
We write $\kappa(x,y,t)$ for the curvature of the curve with coordinates $x$ and $y$ at $t$.
We write $L(x,y,a,b)$ for the length of the curve with coordinates $x$ and $y$ on $[a,b]$.
We write $A(x,y,a,b)$ for the signed area of the curve with coordinates $x$ and $y$ on $[a,b]$.
We write $k_f(x)$ for the curvature of the graph of $f$ at $x$.
We write $\rho_f(x)$ for the radius of curvature of the graph of $f$ at $x$.
We write $P(r,a,b)$ for the area enclosed by the polar curve $r$ on $[a,b]$.
We write $Q(r,a,b)$ for the length of the polar curve $r$ on $[a,b]$.
For every real number $t$, let $c(t)\coloneqq\cos t$.
For every real number $t$, let $s(t)\coloneqq\sin t$.
For every real number $t$, let $p(t)\coloneqq t^2$.
For every real number $t$, let $u(t)\coloneqq t^3$.
For every real number $t$, let $e(t)\coloneqq 1$.

\section{The circle}

\begin{lemma}\label{circle}
Let $t\in\mathbb{R}$. Then $v(c,s,t)=1$ and $\kappa(c,s,t)=1$, and $t$ is a regular point of the curve with coordinates $c$ and $s$.
\end{lemma}

\begin{lemma}\label{length}
$L(c,s,0,2\pi)=2\pi$ and $A(c,s,0,2\pi)=\pi$.
\end{lemma}

\begin{lemma}\label{polar}
$P(e,0,2\pi)=\pi$ and $Q(e,0,2\pi)=2\pi$.
\end{lemma}

\section{Graphs and cusps}

\begin{lemma}\label{parabola}
$k_p(0)=2$ and $\rho_p(0)=\frac12$.
\end{lemma}

\begin{lemma}\label{cusp}
$0$ is a singular point of the curve with coordinates $p$ and $u$.
\end{lemma}

\end{document}
